% =====================================================================
% INTERNAL AI-ASSISTED DRAFT -- NOT AUTHOR-FINALIZED (ICASSP 2027)
% The prose of this manuscript was drafted by an AI coding assistant.
% Experiments, tables, figures, statistics and LaTeX structure are final
% and traceable to generated CSVs (see assets/paper_claims.json), but the
% core prose MUST be substantively rewritten and approved by the human
% author before submission. See HUMAN_REWRITE_CHECKLIST.md.
% Do NOT treat wording, emphasis, framing or title as author-approved.
% Author list / affiliation / contact / ORCID / funding: NOT FILLED.
%
% Figures Fig.~overview and Fig.~structures are COMPILABLE PLACEHOLDERS
% (framed boxes), not final artwork; replace with the drawn figures before
% submission. See FINAL_SUBMISSION_GATE.md.
%
% ICASSP 2027 manuscript. Tables/figures are auto-generated from result CSVs by
% scripts/make_paper_assets.py into assets/; missing assets render as "(pending)"
% so the document always compiles and never contains a hand-typed result.
\documentclass{article}
\usepackage{spconf,amsmath,graphicx,hyperref}
\usepackage{xcolor}
\def\x{{\mathbf x}}
\def\L{{\cal L}}

% ---- include an auto-generated table only if it exists ----
\newcommand{\inputtable}[1]{\IfFileExists{assets/tables/#1}{\input{assets/tables/#1}}{\emph{(table pending: \detokenize{#1})}}}
\newcommand{\figpdf}[2]{\IfFileExists{assets/figs/#1}{\includegraphics[width=#2]{assets/figs/#1}}{\fbox{\parbox{#2}{\centering\vspace{1em}(figure pending: \detokenize{#1})\vspace{1em}}}}}
\setlength{\intextsep}{8pt plus 2pt minus 2pt}
\setlength{\abovecaptionskip}{6pt}
\setlength{\belowcaptionskip}{2pt}
% ---- compilable figure placeholder (not final artwork) ----
\newcommand{\figplaceholder}[3]{\fbox{\parbox[c][#2][c]{0.92\columnwidth}{\centering\footnotesize\itshape
  Figure placeholder\\[2pt]\normalfont\textbf{#1}\\[3pt]\normalfont #3}}}

\title{AUDITING STRUCTURED POLICY LEARNING FOR STREAMING TRANSACTION
COLLATERAL CONTROL: FACTORIZATION, CONDITIONING, AND CROSS-SCALE TRANSFER}

% TODO(authors): confirm full author list, affiliation, and funding before
% submission. Name below is the repository owner; do not infer co-authors.
\name{Yingda Yu\thanks{AI-assisted development disclosure: an AI coding
assistant was used to re-implement code, run experiments, and draft text;
all results are generated by the released code and data, and the author
vets every claim. Author affiliation and funding to be confirmed.}}
\address{Affiliation to be confirmed}

\begin{document}
\ninept
\maketitle

\begin{abstract}
Streaming payment channels must route every arriving transaction into one of
$k$ collateral wallets or drop it, periodically flushing a wallet at a fixed
fee. We study the resulting collateral-control problem as an online,
partially-specified Markov decision process under twelve traffic regimes.
We (i) faithfully port the environment/traffic generator and re-implement three
PPO policies---a joint action head, independent factorized heads, and a
settle-conditioned factored head. Across five seeds, the settle-conditioned
factored head significantly beats the joint head at all four tested
capacities, and the independently factored head at three of four, whereas the
extra settle-conditioning path gives \emph{no statistically detectable} gain
(paired 95\% CIs; zero- and shuffled-conditioning ablations are
indistinguishable), a negative mechanism result we report. (ii) We introduce a
permutation-equivariant set encoder with no $k$-dependent parameter, which
loads and runs zero-shot across wallet counts where a flat MLP cannot; its
transfer quality is asymmetric, and at the largest tested $k$ it trails a
size-matched flat network.
(iii) A deliberately strong, validation-selected threshold rule attains the
highest Money and significantly outperforms the learned policies in nearly all
paired comparisons on stationary streams, and incurs zero avoidable drops
after abrupt regime shifts. On this balanced workload, wallet
fill is close to a sufficient statistic; we keep this negative rather than hide
it. All numbers are produced by released, deterministic code on ported data.
\end{abstract}

\begin{keywords}
reinforcement learning, structured policies, resource allocation,
permutation equivariance, online decision making.
\end{keywords}

\section{Introduction}
Off-chain payment networks such as the Lightning
Network~\cite{poon2016lightning,miller2019sprites,kolachala2024sok} process
streams of transactions over wallets backed by finite collateral; routing under
scarce channel funds has been studied for payment channel
networks~\cite{sivaraman2020spider}.
Transactions arrive sequentially, each carrying a requested value, and at every
arrival the node must decide where to route it: settle it into one of $k$
wallets that currently holds enough free balance, or reject (drop) it.
Independently, it may \emph{flush} one wallet---empty and recharge it at a fixed
fee and with a short cooldown during which the wallet cannot accept funds. This
is an inherently \emph{sequential} online decision: today's flush frees capacity
for tomorrow's arrivals but costs money and briefly removes a wallet from
service, while greedily accepting now can starve a later, larger transaction.
Collateral $C$ is partitioned \emph{per wallet} rather than pooled, so the
controller must jointly reason about \emph{where} value lands and
\emph{when/what} to empty under non-stationary, bursty traffic with no lookahead
to future arrivals. Such sequential decisions are the domain of online
competitive analysis~\cite{borodin1998online} and of its learning-augmented
variants~\cite{lykouris2018competitive}. The closest
formal model, online collateral maintenance with $k$ wallets of capacity $C/k$,
is analyzed with worst-case competitive guarantees
in~\cite{almashaqbeh2024competitive}; related online-inventory and
learning-augmented packing results
include~\cite{lin2019competitive}.

This structure makes naive policy parameterization costly. A joint categorical
controller emits one logit for every (settle, flush) pair, i.e.
$(k{+}1)^2$ outputs at $k=24$, conflating two decisions that are only weakly
coupled and allocating output capacity to invalid or dominated combinations.
Branched and factorized action heads have been proposed for exactly this
reason~\cite{tavakoli2018action}. \emph{Factorization} is therefore the natural
structural move: predict the
settle wallet and the flush wallet as two smaller $k{+}1$-way heads, dropping
the output count to $2(k{+}1)$. One can then ask whether the residual coupling
matters at all---specifically, whether conditioning the flush decision on the
\emph{chosen} settlement wallet (a directed information path
settle$\to$flush) adds value, or whether the two choices are already
independent under the reward. We treat even that plausible mechanism as a
hypothesis to test, not a feature to assume.

Deployment adds a second structural question. The number of wallets $k$ is an
operating condition that can change between deployments, yet a flat multilayer
perceptron bakes $k$ into its weight shapes: a model trained at one $k$ cannot
even be loaded at another. Encoding the wallets as an unordered set with a
shared, permutation-equivariant function removes all $k$-dependent parameters
and allows one trained policy to be deployed zero-shot at a new wallet count.
Whether such portability comes at an accuracy cost at the original scale is
again an empirical question, not a given.

Finally, learned controllers must be audited against a baseline that is
actually strong. The rule baselines used in the prior manuscript are simple
ones (flush everything, or flush when full) that do not use the most informative
available signal---current wallet fill. We instead tune a simple
balance-feedback threshold on validation data and put it on equal footing with
the learned policies. This turns the paper from ``yet another RL policy'' into
an audit of \emph{which structural inductive biases genuinely help online
collateral control}: does factoring the action help, does settle-conditioning
add anything, does an equivariant encoder buy deployability, and can learning
beat a cheap feedback rule that observes the same state? We answer these on a
faithful port of the prior manuscript's environment and twelve-regime generator
with matched seeds and paired confidence intervals, reporting null results as
prominently as positive ones.

\noindent Our contributions are:
\begin{itemize}
\item A faithful, tested re-implementation of the environment and the
twelve-regime traffic generator, and of three PPO policies---\textbf{JA-PPO}
(joint action over $(k{+}1)^2$ outcomes), \textbf{IFAC} (independent settle and
flush heads, $2(k{+}1)$ outputs), and \textbf{SC-FAC} (flush head conditioned
on the chosen settlement wallet). The environment is a faithful port; learned
policies and unspecified hyperparameters are documented re-implementations.
\item A \textbf{permutation-equivariant set encoder} (shared per-wallet
function + pooling + per-wallet logits, in the spirit of Deep
Sets~\cite{zaheer2017deepsets}) with no parameter that depends on $k$, which
transfers \emph{zero-shot} across wallet counts; a flat multilayer perceptron
cannot (its weight shapes change with $k$).
\item \textbf{Mechanism ablations} that zero or shuffle the conditioning
signal over matched seeds; they find no detectable Money benefit from the
settle$\to$flush path at the tested flush price---an honest null result.
\item An \textbf{honest strong baseline}: a validation-selected best-fit
threshold rule that attains zero avoidable drops, is robust to regime shifts,
and attains the highest Money at every capacity, beating the learned policies
in nearly all paired comparisons. We report where learning does---and, more
often, does not---beat it.
\end{itemize}

\section{Streaming collateral control}
\textbf{Model.} Collateral $C\in\{800,900,1000,1200\}$ is split into $k=24$
wallets of capacity $C/k$. At each of $T=1000$ decisions a transaction of size
$x_t\in[0,1000]$ arrives. The agent picks a settle wallet $a_s\in[0,k]$ and a
flush wallet $a_f\in[0,k]$, where index $k$ is a null (no-op). A flush executes
\emph{first}, emptying a wallet and making it unavailable for $F{=}3$ decisions
before it refills to full; the transaction is then accepted if $a_s$ is usable
and has free balance $\ge x_t$. This order is what couples the two choices in a
single step: flushing the wallet into which value is about to settle causes the
transaction to be dropped. Drops are attributed as \emph{oversize}
($x_t>C/k$, physically impossible for any policy), \emph{frozen} (cooldown), or
\emph{avoidable} (insufficient balance or a flush/settle conflict that a
proactive flush could have prevented). Reporting drops by cause lets us separate
the structural floor that no controller can remove from the avoidable fraction
that measures control quality. The observation is $3k{+}2$: per-wallet
normalized balance, availability, and cooldown, plus $x_t$ normalized and a
time-progress scalar (Fig.~\ref{fig:overview}).

\textbf{Objective.} Per-episode Money is
$M=p\!\cdot\!\sum\!\text{accepted}_x-\tau\!\cdot\!\#\text{flushes}$ with
acceptance price $p=1$ and flush fee $\tau=10$, so accepting value and spending
flushes trade off directly in one scalar rather than through separate metrics.
Episodic PPO~\cite{schulman2017ppo} with generalized advantage
estimation~\cite{schulman2015gae} optimizes a stochastic policy; evaluation
uses the deterministic (argmax) policy on fixed pools.

\textbf{Data.} Episodes are drawn from twelve regimes spanning steady and
large-transaction traffic and correlated bursts (calibrated to a target mean),
all generated with fixed seeds: 5000 training, 300 validation, and 200 test
episodes per regime. Every method and seed is evaluated on the \emph{same}
deterministic test pools, so cross-policy comparisons are paired at the episode
level. All hyperparameters absent from the prior manuscript are fixed from
validation pilots, never from the test set; failed seeds are retained rather
than discarded.

\begin{figure}[t]
\centering
\figpdf{overview.pdf}{0.88\columnwidth}
\caption{Streaming collateral control. Each arriving transaction is settled
into one of $k$ capacity-partitioned wallets or dropped; the controller may
additionally flush one wallet at fee $\tau$. Flush executes before settle.}
\label{fig:overview}
\end{figure}

\section{Policies}
\noindent\textbf{JA-PPO} emits one categorical over all settle/flush
pairs: $(k{+}1)^2$ logits ($625$ at $k{=}24$). This treats the step as a single
large joint decision, which both inflates the output layer and forces the
network to allocate probability mass to combinations that are invalid (flushing
and settling the same wallet) or obviously dominated. Structured action
spaces have been addressed with large-action embeddings~\cite{dulac2015deep},
parameterized-action hierarchies~\cite{hausknecht2016deep}, action
branching~\cite{tavakoli2018action}, and neighborhood
construction~\cite{akkerman2024dynamic}.

\noindent\textbf{IFAC} factors the action into two independent $k{+}1$
categoricals (50 logits), $\pi(a_s,a_f\mid s)=\pi_s(a_s\mid s)\,\pi_f(a_f\mid
s)$. This removes the quadratic joint head, encodes the weak-coupling prior
between the two sub-decisions, and shares the trunk representation. It also
makes each sub-decision cheaper to learn: most settle/flush combinations are
never observed, while each head sees its own signal at every step.

\noindent\textbf{SC-FAC} keeps the factorization but routes an embedding of the
\emph{chosen} settle wallet into the flush head,
$\pi_f(a_f\mid s,a_s)$, so flushing can in principle depend on where value was
just placed. The mechanism it is designed to capture is targeted liquidity
management---flushing a wallet other than the one just funded---and is the one
``extra'' inductive bias whose payoff is genuinely uncertain. Two ablations
share its architecture exactly: \textbf{no-cond.} zeroes the settle embedding,
and \textbf{shuffled} feeds a deranged (wrong-wallet) embedding. Comparing the
three under matched seeds isolates whether the information path carries any
causal value.

\noindent\textbf{Set encoder.} The $k$ wallets are intrinsically an unordered
set: permuting wallet identities must not change the decision. A flat MLP
violates this implicitly by treating wallet slots as an ordered, fixed-length
vector, and---more importantly for deployment---ties every hidden weight to the
input width and hence to $k$. We instead encode each wallet with a shared
function $h_i=\phi(\text{balance}_i,\allowbreak\text{avail}_i,\allowbreak\text{cooldown}_i,g)$ of its
state and two global scalars, pool $\{h_i\}$ symmetrically, and produce
per-wallet settle/flush logits (Fig.~\ref{fig:structures}). The map is
permutation-equivariant by construction (verified by unit tests), and no layer
size depends on $k$; this \emph{$k$-independence} is exactly the property a
flat network lacks, and it is what lets weights trained at one wallet count be
loaded and executed---without retraining---at another. The settle-conditioned
set variant gathers the chosen settle wallet's feature into the flush head.

\begin{figure}[t]
\centering
\figpdf{structures.pdf}{0.82\columnwidth}
\caption{Action parameterizations. JA-PPO uses a quadratic joint head; IFAC
uses two linear factorized heads; SC-FAC adds a settle-conditioned edge to the
flush head; the set variant replaces $k$-shaped trunk weights with a shared
equivariant encoder, enabling zero-shot cross-$k$ deployment.}
\label{fig:structures}
\end{figure}

\begin{table}[htb]
\centering
\footnotesize
\inputtable{efficiency.tex}
\caption{Parameters, output logits, and CPU latency at $C{=}1200,\ k{=}24$.
Set-encoder params are independent of $k$ (transferable).}
\label{tab:eff}
\end{table}

\section{Experiments}
We evaluate deterministic policies on fixed, shared test pools. Model and
training choices are made on validation only; failed seeds and null results are
retained. \emph{Baselines} are the native one-flush references FA and FWF
(reconstructed, since the reference design used multiple flushes per step), a
rotation rule, and the strong \textbf{BFP0.5}: best-fit routing that
proactively flushes the fullest usable wallet once its fill crosses a
threshold $0.5$, chosen on validation. We also compute an oracle feasible value
(transactions that fit at least one wallet) and decompose drops by cause.
Results are organized around five questions: whether factorization beats the
joint head; whether conditioning adds anything; whether equivariance buys
cross-$k$ transfer and at what cost; who is reliable under abrupt regime
shifts; and whether the structural changes improve parameter/output efficiency.

\noindent\textbf{Stationary performance: does factorization help?}
Table~\ref{tab:main} reports Money per capacity and
Table~\ref{tab:paired} the seed-level paired contrasts at $C{=}1200$; the
structural oversize floor is $\approx400$ drops per episode for \emph{every}
method, so the differences live in the avoidable fraction. Two robust effects
and one null emerge. \emph{Factorization helps:} SC-FAC significantly
outperforms JA-PPO at all four capacities ($p\!\le\!0.013$, $5/5$ seeds), and
IFAC outperforms JA-PPO at $C\!\in\!\{800,900,1200\}$ ($p\!\le\!0.033$) but
\emph{not} at $C{=}1000$ ($p{=}0.57$). The mechanism interpretation is direct:
avoidable drops fall from $77$/episode (JA-PPO) to $49.5$ (IFAC) and $40.4$
(SC-FAC), exactly the failure mode the smaller, better-coupled heads should
remove. \emph{Conditioning does not add Money:} the SC-FAC$-$IFAC paired
difference has a 95\% CI including zero at $C\!\in\!\{900,1000,1200\}$, and at
$C{=}800$ SC-FAC is significantly \emph{lower} than IFAC ($-419$, $p{=}0.009$);
SC-FAC accepts marginally more value but spends more flushes, which cancels the
gain at $\tau{=}10$. \emph{A tuned rule beats learning:} BFP0.5 attains the
highest Money at every capacity with zero avoidable drops, and its paired
advantage is significant in $11/12$ method-capacity contrasts; the sole
exception is vs.\ IFAC at $C{=}1000$ ($+2545$, $p{=}0.057$, CI crossing zero),
which BFP0.5 nevertheless wins on all five seeds. On this balanced i.i.d.\
workload a single balance-feedback threshold is near-optimal, and a
3000-episode PPO budget does not beat it (Fig.~\ref{fig:cap}).

\begin{table}[htb]
\centering
\scriptsize
\setlength{\tabcolsep}{3.2pt}
\inputtable{main_results.tex}
\caption{Money per episode (mean$\pm$SE over seeds) vs.\ collateral $C$.
Rules need no seeds; prior-manuscript reported values are not shown (they
remain in the reproduction report).}
\label{tab:main}
\end{table}

\begin{table}[htb]
\centering
\scriptsize
\setlength{\tabcolsep}{3.2pt}
\inputtable{paired.tex}
\caption{Seed-level paired Money differences at $C{=}1200$ (mean $a-b$,
paired-$t$ 95\% CI, $p$-value, wins over seeds). Positive favors $a$.}
\label{tab:paired}
\end{table}

\begin{figure}[htb]
\centering
\figpdf{capacity_curve.pdf}{0.40\columnwidth}
\caption{Money vs.\ collateral $C$ (solid: learned; dashed: rules). BFP0.5
dominates at every capacity; the factorized policies separate from JA-PPO.}
\label{fig:cap}
\end{figure}

\noindent\textbf{Does settle-conditioning carry causal value?}
Table~\ref{tab:abl} compares SC-FAC against the no-conditioning and
shuffled-conditioning variants under matched parameters, seeds, and budget.
Neither zeroing nor permuting the settle embedding lowers Money detectably:
every paired 95\% CI includes zero, and at $C{=}1200$ the fully conditioned
variant is nominally the lowest of the three. A correct but useless wiring is
not evidence against architecture---it is evidence that, at the tested flush
price and observation model, knowing the chosen settle wallet is not
decision-relevant once wallet state is already observed.

\begin{table}[htb]
\centering
\footnotesize
\inputtable{ablation.tex}
\caption{Settle-conditioning mechanism ablations ($C{=}1200$).}
\label{tab:abl}
\end{table}

\noindent\textbf{Does equivariance buy cross-$k$ transfer?}
We train at $k\in\{6,12,24\}$ and deploy, with no retraining, at every $k$.
Fig.~\ref{fig:heat} shows that the set encoder loads and runs at every
deployment $k$, whereas flat policies have no off-diagonal entry at all because
their state-dictionary shapes change with $k$. The value of the structure is
\emph{operational deployability}, not uniformly preserved Money: zero-shot
retention is asymmetric---strong when moving to larger-wallet/smaller-$k$
operating points and weak when a model trained with large wallets meets the
small wallets at $k{=}24$ (e.g.\ $k{=}6{\to}24$ yields $4433$). At matched $k$
the set encoder is comparable to or above a flat MLP at $k{=}6,12$ but lower at
$k{=}24$ ($9190$ vs.\ $13610$): portability comes with a scale-specific accuracy
trade-off, concentrated at the hardest (many small wallets) operating point.

\begin{figure}[htb]
\centering
\figpdf{transfer_set_sc_fac.pdf}{0.40\columnwidth}
\caption{Cross-$k$ transfer matrix (Set-SC-FAC): rows train $k$, columns
deploy $k$. Off-diagonal cells are zero-shot (flat MLP has none).}
\label{fig:heat}
\end{figure}

\begin{sloppypar}
\noindent\textbf{Who is reliable under regime shifts?}
Trained only on stationary pools, every policy is tested on the \emph{same}
deterministic episodes spliced across regimes at an abrupt change point (six
scenarios~\cite{gama2014survey}: calm$\to$burst, burst$\to$calm, early and
late).
Table~\ref{tab:sw} reports post-switch avoidable drops and Money. BFP0.5 incurs
\emph{zero} avoidable drops after every switch and the highest Money, adapting
without retraining because current fill is close to a sufficient statistic for
this feedback law. The learned policies generalize far better than the naive
rules ($28$--$56$ vs.\ $144$--$148$ post-switch drops). Against BFP0.5, JA-PPO
and IFAC remain worse on both drops and Money over the paired scenarios
(one-sample paired tests, $n{=}3$; $p\!<\!0.01$); for SC-FAC the drops
contrast is weak ($p\!\approx\!0.09$), and its Money contrast
($p\!<\!0.05$ uncorrected) does not survive Holm correction across
the six switching tests ($p\!\approx\!0.07$). IFAC and SC-FAC incur similar post-switch drops ($28.3$). The defensible reading is that learned policies are
robust enough not to collapse when the stream changes---but we do
\emph{not} observe any switch-adaptation advantage of learning over the tuned
feedback rule, and we deliberately do not claim one.

\begin{table}[htb]
\centering
\footnotesize
\setlength{\tabcolsep}{4pt}
\inputtable{switching.tex}
\caption{Switching streams (mean over six scenarios; learned over three seeds).
Post-switch avoidable drops and Money.}
\label{tab:sw}
\end{table}
\end{sloppypar}

\noindent\textbf{Parameter and output efficiency.}
Table~\ref{tab:eff} shows the structural payoff independently of Money.
Replacing the joint head cuts outputs from $625$ to $50$ and the actor from
$245{,}874$ to $98{,}099$ parameters (IFAC); the set variants use $135$k
parameters with \emph{no} $k$-dependent tensor and sub-millisecond-per-step CPU
latency throughout. Factorization therefore yields a representation that is
smaller and cheaper at the same time as it is statistically stronger than the
joint head on most capacities.

\section{Discussion and limitations}
The central result is a boundary map for structured policy learning on this
task, not a blanket win for RL. \emph{Factorization is the reliable learned
effect}: replacing a quadratic joint head with two linear sub-heads removes
exactly the avoidable drops the joint head was producing, with gains that
survive multi-seed paired tests at nearly every capacity and a simultaneous
reduction in parameters and outputs. \emph{Settle-conditioning is a null}:
once wallet fill is in the state, the identity of the just-chosen settle
wallet carries no additional decision-relevant signal at $\tau{=}10$, and both
the zero and shuffled controls confirm this. We report it as a wiring that is
plausible but unremunerative here, rather than quietly dropping it.

\emph{Equivariance buys a different currency.} The set encoder's value is not
matched-$k$ accuracy---it is the ability to ship one trained policy to a
deployment whose wallet count the training never saw, which a flat MLP
structurally cannot do. Its weaknesses are honest: transfer is asymmetric, and
at the hardest point ($k{=}24$ small wallets) the inductive bias costs
matched-scale accuracy. A practitioner operating one fixed $k$ should prefer
the flat factorized model; one who must redeploy across changing $k$ gains
real optionality.

\emph{The strong rule is the paper's uncomfortable finding.} On balanced,
i.i.d.\ streams with full state observation, a one-line fill-feedback
threshold reaches zero avoidable drops stationary and after every abrupt
switch, and beats the learned controllers in nearly every paired comparison.
This must not be hidden or reframed: it means that on the current observation
model, wallet fill is close to a sufficient statistic and the residual problem
is nearly solved by feedback control. It also defines precisely where future
learning could earn its place---predictable non-i.i.d.\ burst structure that
fill alone cannot anticipate, heterogeneous wallet sizes or fees that break the
single-threshold policy, multi-step lookahead under expensive flushes, or
reward functions deliberately misaligned with accepted value (a candidate
explanation for SC-FAC's extra flushing that future work could test). Better
exploration and vectorized, larger-budget PPO training are engineering levers
toward the same gap; we deliberately do not instantiate an untested
heterogeneous environment here to manufacture a regime where RL wins.

Limitations are concrete: one flush per step, finite horizon, a single fixed
fee, no lookahead to future arrivals, homogeneous wallets, and single-process
PPO without vectorized rollouts; CPU training was also faster than the available
GPU for these batch-size-one rollouts. The prior manuscript's learned numbers
are treated as reported, not regenerated: the environment is a faithful port
while the policies and unspecified hyperparameters are documented
re-implementations chosen on validation.

\section{Conclusion}
We audited which structural inductive biases help a learned streaming
collateral controller. The evidence is positive on two points, null on one, and
sobering on one. Factoring the joint action into separate settle/flush heads
gives a reproducible gain over the joint head---significant at all four
capacities for the settle-conditioned variant and at three of four for the
independent variant---and cuts outputs from $(k{+}1)^2$ to $2(k{+}1)$. Adding a
settle-conditioned flush path yields no detectable Money gain under
matched-seed zero/shuffle ablations. A $k$-independent equivariant encoder
adds genuine operational value---weights deploy zero-shot to unseen wallet
counts where a flat MLP cannot---at an asymmetric transfer and a matched-scale
cost concentrated at the hardest operating point. Most importantly, on
homogeneous stationary streams and across abrupt regime switches, one
validation-selected feedback rule attains the highest Money with zero
post-switch avoidable drops and beats the learned controllers in nearly every
paired comparison. We present that rule as the benchmark future learning
methods must beat, and the audit---what helps, what does not---as the
contribution.

\clearpage
\bibliographystyle{IEEEbib}
\bibliography{refs}
\end{document}
